\documentclass[10pt, letterpaper]{article}

% Packages:
\usepackage[
ignoreheadfoot, % set margins without considering header and footer
top=1.2 cm, % seperation between body and page edge from the top
bottom=1.2 cm, % seperation between body and page edge from the bottom
left=1.5 cm, % seperation between body and page edge from the left
right=1.5 cm, % seperation between body and page edge from the right
footskip=1.0 cm, % seperation between body and footer
% showframe % for debugging 
]{geometry} % for adjusting page geometry
\usepackage{titlesec} % for customizing section titles
\usepackage{tabularx} % for making tables with fixed width columns
\usepackage{array} % tabularx requires this
\usepackage[dvipsnames]{xcolor} % for coloring text
\definecolor{primaryColor}{RGB}{0, 79, 144} % define primary color
\usepackage{enumitem} % for customizing lists
\usepackage{fontawesome5} % for using icons
\usepackage{amsmath} % for math
\usepackage[
pdftitle={Hamid Abbasi's CV},
pdfauthor={Hamid Abbasi},
pdfcreator={LaTeX with RenderCV},
colorlinks=true,
urlcolor=primaryColor
]{hyperref} % for links, metadata and bookmarks
\usepackage[pscoord]{eso-pic} % for floating text on the page
\usepackage{calc} % for calculating lengths
\usepackage{bookmark} % for bookmarks
\usepackage[symbol]{footmisc}
\usepackage{lastpage} % for getting the total number of pages
\usepackage{changepage} % for one column entries (adjustwidth environment)
\usepackage{lipsum}  % for dummy text
\usepackage{paracol} % for two and three column entries
\usepackage{ifthen} % for conditional statements
\usepackage{needspace} % for avoiding page brake right after the section title
\usepackage{iftex} % check if engine is pdflatex, xetex or luatex

% Ensure that generate pdf is machine readable/ATS parsable:
\ifPDFTeX
\input{glyphtounicode}
\pdfgentounicode=1
% \usepackage[T1]{fontenc} % this breaks sb2nov
\usepackage[utf8]{inputenc}
\usepackage{lmodern}
\fi



% Some settings:
\AtBeginEnvironment{adjustwidth}{\partopsep0pt} % remove space before adjustwidth environment
\pagestyle{empty} % no header or footer
\setcounter{secnumdepth}{0} % no section numbering
\setlength{\parindent}{0pt} % no indentation
\setlength{\topskip}{0pt} % no top skip
\setlength{\columnsep}{0cm} % set column seperation
\makeatletter
\let\ps@customFooterStyle\ps@plain % Copy the plain style to customFooterStyle
\patchcmd{\ps@customFooterStyle}{\thepage}{
\color{gray}\textit{\small Hamid Abbasi - Page \thepage{} of \pageref*{LastPage}}
}{}{} % replace number by desired string
\makeatother
\pagestyle{customFooterStyle}

\titleformat{\section}{\needspace{4\baselineskip}\bfseries\large}{}{0pt}{}[\vspace{1pt}\titlerule]

\titlespacing{\section}{
% left space:
-1pt
}{
% top space:
0.3 cm
}{
% bottom space:
0.2 cm
} % section title spacing

\renewcommand\labelitemi{$\circ$} % custom bullet points
\newenvironment{highlights}{
\begin{itemize}[
topsep=0.10 cm,
parsep=0.10 cm,
partopsep=0pt,
itemsep=0pt,
leftmargin=0.4 cm + 10pt
]
}{
\end{itemize}
} % new environment for highlights

\newenvironment{highlightsforbulletentries}{
\begin{itemize}[
topsep=0.10 cm,
parsep=0.10 cm,
partopsep=0pt,
itemsep=0pt,
leftmargin=10pt
]
}{
\end{itemize}
} % new environment for highlights for bullet entries


\newenvironment{onecolentry}{
\begin{adjustwidth}{
0.2 cm + 0.00001 cm
}{
0.2 cm + 0.00001 cm
}
}{
\end{adjustwidth}
} % new environment for one column entries

\newenvironment{twocolentry}[2][]{
\onecolentry
\def\secondColumn{#2}
\setcolumnwidth{\fill, 4.5 cm}
\begin{paracol}{2}
}{
\switchcolumn \raggedleft \secondColumn
\end{paracol}
\endonecolentry
} % new environment for two column entries

\newenvironment{header}{
\setlength{\topsep}{0pt}\par\kern\topsep\centering\linespread{1.5}
}{
\par\kern\topsep
} % new environment for the header

\newcommand{\placelastupdatedtext}{% \placetextbox{<horizontal pos>}{<vertical pos>}{<stuff>}
\AddToShipoutPictureFG*{% Add <stuff> to current page foreground
\put(
\LenToUnit{\paperwidth-2 cm-0.2 cm+0.05cm},
\LenToUnit{\paperheight-1.0 cm}
){\vtop{{\null}\makebox[0pt][c]{
\small\color{gray}\textit{Last updated in December 2024}\hspace{\widthof{Last updated in September 2024}}
}}}%
}%
}%

% save the original href command in a new command:
\let\hrefWithoutArrow\href

% new command for external links:
\renewcommand{\href}[2]{\hrefWithoutArrow{#1}{\ifthenelse{\equal{#2}{}}{ }{#2 }\raisebox{.15ex}{\footnotesize \faExternalLink*}}}


\begin{document}
\newcommand{\AND}{\unskip
\cleaders\copy\ANDbox\hskip\wd\ANDbox
\ignorespaces
}
\newsavebox\ANDbox
\sbox\ANDbox{}

\placelastupdatedtext
\begin{header}
\textbf{\fontsize{24 pt}{24 pt}\selectfont Hamid Abbasi}

\vspace{0.3 cm}

\normalsize
% \mbox{{\color{black}\footnotesize\faMapMarker*}\hspace*{0.13cm}Manchester, UK}%
% \kern 0.25 cm%
% \AND%
% \kern 0.25 cm%
\mbox{\hrefWithoutArrow{mailto:h.abbasi@imperial.ac.uk}{\color{black}{\footnotesize\faEnvelope[regular]}\hspace*{0.13cm}h.abbasi@imperial.ac.uk}}%
\kern 0.25 cm%
\AND%
\kern 0.25 cm%
\mbox{\hrefWithoutArrow{tel:+44-7575-123-808}{\color{black}{\footnotesize\faPhone*}\hspace*{0.13cm}07575 123 808}}%
% \kern 0.25 cm%
% \AND%
% \kern 0.25 cm%
% \mbox{\hrefWithoutArrow{https://scholar.google.com/citations?user=uMCSKucAAAAJ&hl=en}{\color{black}{\footnotesize\faGraduationCap}\hspace*{0.13cm}Google Scholar}}%
\kern 0.25 cm%
\AND%
\kern 0.25 cm%
\mbox{\hrefWithoutArrow{https://www.linkedin.com/in/hamid-abbasi-0a997114a/}{\color{black}{\footnotesize\faLinkedinIn}\hspace*{0.13cm}Hamid Abbasi}}%
\kern 0.25 cm%
\AND%
\kern 0.25 cm%
\mbox{\hrefWithoutArrow{https://github.com/HamidAbbasi-R}{\color{black}{\footnotesize\faGithub}\hspace*{0.13cm}HamidAbbasi-R}}%
\kern 0.25 cm
\AND
\kern 0.25 cm
\mbox{\color{black}{\footnotesize\faMap}\hspace*{0.13cm}London, UK}%
\end{header}

\section{Work Experience}

\begin{twocolentry}{\textit{July 2025 – Present}}
\textbf{Research Associate}, London, Imperial College London
\end{twocolentry}

\begin{highlights}
\item Project \textbf{CASTLE} (Contrail Analysis using Stochasticity and Thermodynamic consistency in a Lagrangian Environment): Architected a scalable Lagrangian model of contrail-cirrus atmospheric interactions in C++ and Python, leveraging hybrid parallelization (MPI/OpenMP) on HPC cluster across distributed nodes; addressing the physically convoluted and computationally expensive problem of regional-scale contrail modelling. \textit{Github repo currently private pending publication; source code available upon request.}
\item Developed a novel stochastic model for cirrus cloud formation, explicitly linked to probabilistic parameterizations of atmospheric gravity waves to achieve an accurate simulation of the atmospheric ice life-cycle. 
\item Project \textbf{CLAMP}: Optimized legacy Python code by orders of magnitude, reducing execution time through advanced vectorization and Numba compilation; creating a novel and highly efficient plume-scale representation of contrail dynamics. \textit{Github repo currently private pending publication; source code available upon request.}
\end{highlights}

\begin{twocolentry}{\textit{Sep 2022 – Jan 2024}}
\textbf{Graduate Teaching Assistant}, Manchester, University of Manchester
\end{twocolentry}

\begin{highlights}
\item Process Design: Teaching and helping students with construction of P\&ID designs
\item Engineering Math 2 \& 3: Linear algebra, ODEs \& PDEs, Multi variable Calculus, Complex Analysis
\item Fundamentals of Thermodynamics: Analysis of 1\textsuperscript{st} and 2\textsuperscript{nd} laws in various thermal systems.
\end{highlights}



\begin{twocolentry}{\textit{Sep 2023 – Jan 2024}}
\textbf{Lecturer}, Manchester, Wanhua Chemical Group
\end{twocolentry}

\begin{highlights}
% \item Reaction Systems: Covered key topics such as reactor performance, chemical equilibrium, idealized reactor models, stochastic simulations of chemical reactions (Gillespie algorithm), and multiphase reactors. Delivered in-depth instruction on the principles of reaction engineering.
\item Reaction Systems: Conducted sessions on the application of Monte Carlo methods (Gillespie algorithm), in the context of simulating the stochastic time-dependent dynamics of chemical reactions.
\end{highlights}


\begin{twocolentry}{\textit{Sep 2018 – Jan 2021}}
\textbf{Graduate Teaching Assistant}, Iran University of Science and Technology
\end{twocolentry}

\begin{highlights}
\item Optimization of Thermal Systems; Thermodynamics and CFD; Advanced IC Engines; Fluid Mechanics 1 \& 2
\end{highlights}


\section{Education}

\begin{twocolentry}{

\textit{Jan 2022 – Jul 2025}}
\textbf{University of Manchester}\textit{, PhD in Chemical Engineering}
\end{twocolentry}

% \vspace{0.10 cm}
\begin{highlights}
\item \textbf{Thesis:} multiscale model of solid oxide fuel cells \href{https://github.com/HamidAbbasi-R/SOFC-project}{\faGithub\hspace{0.13cm}/SOFC-Project}
\item Awarded the University of Manchester’s flagship scholarship (\href{https://www.manchester.ac.uk/study/postgraduate-research/funding/presidents-doctoral-scholar-award/}{President's Doctoral Scholarship})
\item Designed a parallelized finite-volume model in Python to model and solve nonlinear PDEs at different length scales.
% \item Developed a numerical framework to generate large synthetic 3D images mimicking porous microstructure of electrodes, and applied statistical concepts to extract features and make comparisons with the real materials.
\item Developed an integrated pore-scale, and cell-scale model to simulate chemical reactions involved in the internal reforming of hydrocarbon fuels within porous electrodes of SOFC accounting for chemical non-equilibrium.
% \item Designed custom Python pipelines for 3D interactive visualization and analysis of tomographic datasets, addressing unique research-specific challenges such as triple-phase boundary detection and visualization.
\end{highlights}


\begin{twocolentry}{
\textit{Sep 2017 – June 2020}}
\textbf{Iran University of Science and Technology}, \textit{MSc in Mechanical Engineering}
\end{twocolentry}

% \vspace{0.10 cm}
\begin{highlights}
\item GPA: 3.88/4.0 (\href{https://drive.google.com/file/d/1WRgpXZXEvMWSUbna5VOuGOtFu3cHqanD/view?usp=drive_link}{Official transcript}); Ranked 1\textsuperscript{st} among 2020 graduates.
\item \textbf{Thesis:} Developed a numerical model to solve PDEs governing transient turbulent airflow inside air compressors.
\end{highlights}


\begin{twocolentry}{
\textit{Sep 2013 – June 2017}}
\textbf{Shiraz University}\textit{, BSc in Mechanical Engineering}
\end{twocolentry}

% \vspace{0.10 cm}
\begin{highlights}
\item GPA: 3.55/4.0 (\href{https://drive.google.com/file/d/14eRQ84CCNWGHTbv2qsUlqq4ChiFKmclX/view?usp=drive_link}{Official transcript}) Among top 10\% of graduating students. 
\end{highlights}


\section{Selected Publications}
\textbf{Full list: \href{https://scholar.google.com/citations?user=uMCSKucAAAAJ&hl=en}{Google Scholar}; h-index=10, Citations=829;} Published 16 articles in leading journals, 11 first-author and 6 conferences; International collaborations with EPFL Switzerland, UBC Canada, UNSW Sydney, University of Leeds

% \vspace{-0.2 cm}
 
\begin{highlights}
\item A Pore-Scale Computational Framework for Enhancing Solid Oxide Fuel Cell Performance through the Design of Anode Porous Configurations; 2025, \textit{\textbf{H Abbasi}, et al.} \href{https://doi.org/10.1149/1945-7111/adbfc4}{10.1149/1945-7111/adbfc4}
\item Multiscale Modeling of Solid Oxide Fuel Cells Using Various Microscale Domains Across the Length of the Cell; 2024, \textit{\textbf{H Abbasi}, et al.} \href{https://doi.org/10.1016/B978-0-443-28824-1.50387-2}{10.1016/B978-0-443-28824-1.50387-2}
% \item Multi-scale model of solid oxide fuel cell: enabling microscopic solvers to predict physical features over a macroscopic domain; 2023, \textit{\textbf{H Abbasi}, et al.} \href{https://doi.org/10.1016/B978-0-443-15274-0.50166-9}{10.1016/B978-0-443-15274-0.50166-9}
% \item Multi-scale model of solid oxide fuel cell: enabling microscopic solvers to predict physical features over a macroscopic domain; \textit{\textbf{HR Abbasi}, et al.} \href{https://doi.org/10.1016/B978-0-443-15274-0.50166-9}{10.1016/B978-0-443-15274-0.50166-9}
% \item Optimal design of a hybrid thermal-and membrane-based desalination unit based on renewable geothermal energy; 2021, \textit{MH Mohammadi, \textbf{H Abbasi}, MH M Ghodrat.} \href{https://doi.org/10.1016/j.ecmx.2021.100124}{10.1016/j.ecmx.2021.100124}
\item Application of machine learning tools for constrained multi-objective optimization of an HCCI engine; 2021 \textit{A Gharehghani, \textbf{H Abbasi}, P Alizadeh. } \href{https://doi.org/10.1016/j.energy.2021.121106}{10.1016/j.energy.2021.121106}  
% \item Mixed convection enhancement by using optimized porous media and nanofluid in a cavity with two rotating cylinders; \textit{M Tahmasbi, M Siavashi, \textbf{HR Abbasi}, M Akhlaghi.} \href{https://doi.org/10.1007/s10973-020-09604-z}{10.1007/s10973-020-09604-z}
\item Thermal optimization of shell and tube heat exchanger using porous baffles; 2020, \textit{MH Mohammadi, \textbf{H Abbasi}, et al.} \href{https://doi.org/10.1016/j.applthermaleng.2020.115005}{10.1016/j.applthermaleng.2020.115005}
\item Shape optimization of segmental porous baffles for enhanced thermo-hydraulic performance of shell-and-tube heat exchanger; 2020, \textit{\textbf{H Abbasi}, et al.} \href{https://doi.org/10.1016/j.applthermaleng.2020.115835}{10.1016/j.applthermaleng.2020.115835}
% \item Multi-criteria optimization of a renewable hydrogen and freshwater production system using HDH desalination unit and thermoelectric generator; 2020, \textit{\textbf{HR Abbasi}, H Pourrahmanni.} \href{https://doi.org/10.1016/j.enconman.2020.112903}{10.1016/j.enconman.2020.112903}
% \item Multi-objective optimization and exergoeconomic analysis of a continuous solar-driven system with PCM for power, cooling and freshwater production; 2020, \textit{\textbf{H Abbasi}, et al.} \href{https://doi.org/10.1016/j.enconman.2020.112761}{10.1016/j.enconman.2020.112761}
% \item Exergoeconomic optimization of a solar driven system with reverse osmosis desalination unit and phase change material thermal energy storages; \textit{\textbf{HR Abbasi}, et al.} \href{https://doi.org/10.1016/j.enconman.2019.112042}{10.1016/j.enconman.2019.112042}
\end{highlights}

% \textit{Waste heat management of direct carbon fuel cell with advanced supercritical carbon dioxide power cycle – A thermodynamic-electrochemical modeling approach} \textbf{HR Abbasi}, A Yavarinasab, S Roohbakhsh. \href{https://doi.org/10.1016/j.jcou.2021.101630}{10.1016/j.jcou.2021.101630}

% \item \textit{Thermodynamic analysis of a tri-generation system using SOFC and HDH desalination unit} \textbf{HR Abbasi}, H Pourrahmani, N Chitgar, J Van herle. \href{https://doi.org/10.1016/j.ijhydene.2021.04.152}{10.1016/j.ijhydene.2021.04.152}


\section{Certificates}

% \clearpage      % delete if you modified the structure later on
\begin{onecolentry}
\textbf{nVIDIA:} Accelerating CUDA C++ Applications with Multiple GPUs (\href{https://learn.nvidia.com/certificates?id=XDOtXOsrRcWrLlQKnlnjgw}{Link}), Fundamentals of Accelerated Computing with CUDA Python (\href{https://learn.nvidia.com/certificates?id=3xKR2CZzTTqemX0KZ4GvKQ}{Link}) 
\end{onecolentry}
\begin{highlights}
\item Utilized all available GPUs on a single node to scale computation-intensive workloads.
\item Gained experience in GPU-accelerated High Performance Computing (HPC) with CUDA Python using Numba, covering custom kernels, memory coalescing, shared memory, multi-dimensional grids.
% \item Used Nsight Systems Visual Profiler timeline to observe improvement opportunities to achieve peak performance.
\end{highlights}

\begin{onecolentry}
\textbf{nVIDIA:} Deep Learning (\href{https://learn.nvidia.com/certificates?id=kt_dwUnHRfCzbJKPbQE-rw}{Link}), and Transformer-Based Natural Language Processing (\href{https://learn.nvidia.com/certificates?id=uh-SlSjNQ9unOu0fuISakA}{Link}) 
\end{onecolentry}
\begin{highlights}
\item Built, feature-engineered, and fine-tuned deep learning models with PyTorch, covering CNN, efficient transfer learning techniques (from pre-trained models), data augmentation, and natural language processing.
% \item Hands-on experience with transformer-based large language models (LLMs) for advanced NLP tasks, including text classification, named-entity recognition, author attribution, and question answering.
\end{highlights}

\begin{onecolentry}
\textbf{nVIDIA:} Accelerating End-to-End Data Science Workflows (\href{https://learn.nvidia.com/certificates?id=sGjX3cWXSUyOxIFzh7y0Gg}{Link})
\end{onecolentry}
\begin{highlights}
\item Used GPU-accelerated framework RAPIDS (cuDF, cuML, dask) to explore causal relationships in healthcare data, and create interactive visualizations from large datasets uncovering insights in complex systems.
% \item Used GPU-accelerated frameworks (RAPIDS: cuDF, cuML, Dask) to analyze large-scale datasets, including UK census data, GB road networks, and NHS-derived epidemic models, enabling scalable probabilistic modeling.
% \item Explored causal relationships in healthcare data (e.g., Ebola risk factors) through statistical analysis and interactive visualization dashboards, uncovering insights in complex systems.
\end{highlights}


% \begin{onecolentry}
% \textbf{nVIDIA:} Introduction to Transformer-Based Natural Language Processing (\href{https://learn.nvidia.com/certificates?id=uh-SlSjNQ9unOu0fuISakA}{Link})
% \end{onecolentry}
% \begin{highlights}
% \item Hands-on experience with transformer-based large language models (LLMs) for advanced NLP tasks, including text classification, named-entity recognition, author attribution, and question answering.
% \end{highlights}


\begin{onecolentry}
\textbf{HarvardX:} Data Science: Probability, Inference and Modeling, Data Wrangling, Linear Regression, Machine Learning
\end{onecolentry}
\begin{highlights}
\item Applied statistical inference, linear regression, and probability theory to solve real-world challenges. Gained hands-on experience in confidence intervals, Bayesian modeling, and predictive analytics through practical projects.
\item Developed interactive dashboards to explore multiple real-world problems using advanced statistics and probability. See \href{https://github.com/HamidAbbasi-R/Poll-Simulation}{\faGithub\hspace{0.13cm}/Poll-Simulation}, \href{https://github.com/HamidAbbasi-R/Loan-Risk-Analysis}{\faGithub\hspace{0.13cm}/Loan-Risk-Analysis}, and \href{https://github.com/HamidAbbasi-R/Bayesian-Inference}{\faGithub\hspace{0.13cm}/Bayesian-Inference}
\end{highlights}


\begin{onecolentry}
\textbf{Kaggle:} \href{https://www.kaggle.com/learn/certification/hamidrabbasi/intro-to-deep-learning}{Intro} to Deep Learning, \href{https://www.kaggle.com/learn/certification/hamidrabbasi/intro-to-machine-learning}{Intro} and  \href{https://www.kaggle.com/learn/certification/hamidrabbasi/intermediate-machine-learning}{Intermediate} Machine Learning, and \href{https://www.kaggle.com/learn/certification/hamidrabbasi/time-series}{Time Series}
\end{onecolentry}
\begin{highlights}
\item Built and optimized DNNs and ensemble models (Random Forest, XGBoost) on real-world datasets, applying techniques like dropout, batch normalization, and cross-validation to prevent overfitting and ensure robust pipelines. 
% \item Time series analysis; combining linear regression and ensemble learning to capture trends, seasonality, and cycles.
\end{highlights}

\begin{onecolentry}
\textbf{University of Manchester at Harwell:} Visualisation and Quantification of Tomographic Datasets
\end{onecolentry}
\begin{highlights}
\item Hands-on experience with tomographic data analysis, combining theoretical insights and software experience.
\end{highlights}


\section{Selected Projects}

\begin{twocolentry}{\textit{2025-2026}}
\textbf{Contrail/Atmospheric Models}
\end{twocolentry}
\begin{highlights}
\item \textbf{Project CASTLE:} Developed a highly scalable, multi-scale Lagrangian framework for contrail modeling, designed for seamless integration with a diverse range of plume-scale contrail sub-models.
\item \textbf{Project CLAMP:} Developed a high-performance Lagrangian model of plume-scale contrail dynamics.
\end{highlights}

\begin{twocolentry}{\textit{2024}}
\textbf{Statistics and Probabilistic Models}
\end{twocolentry}
\begin{highlights}
\item \href{https://github.com/HamidAbbasi-R/Physics-Informed-ML}{\faGithub\hspace{0.13cm}/Physics-Informed-ML}: Created an interactive tool using PyTorch to model nonlinear dynamics of Van der Pol oscillator using a Tensor Basis Neural Network (TBNN), allowing real-time training based on user inputs.
\item \href{https://github.com/HamidAbbasi-R/HMM-synthetic}{\faGithub\hspace{0.13cm}/HMM-synthetic}: Designed an interactive dashboard to generate synthetic time series data with specific Markov properties, leveraging Hidden Markov Models to evaluate and analyze its performance on these simulated datasets.
\item \href{https://github.com/HamidAbbasi-R/Network-Autoregressive-Model}{\faGithub\hspace{0.13cm}/Network-Autoregressive-Model}: Developed a Python-based tool to simulate and estimate discrete-time linear autoregressive models for networked systems, incorporating adjacency constraints, external inputs, and noise.
\item \href{https://github.com/HamidAbbasi-R/ML-Toolbox-Finance}{\faGithub\hspace{0.13cm}/ML-Toolbox-Finance}: Developed efficient pipelines in Python to quickly evaluate the performance of different classes of ML models (Deep Learning, Hidden Markov Model, Random Forest) forecasting an asset's performance.
\item \href{https://github.com/HamidAbbasi-R/Markov-Chain}{\faGithub\hspace{0.13cm}/Markov-Chain} : Built an interactive dashboard in Python to visualize Markov transition matrices of assets.
% \item Trained the TBNN with a combination of data-driven and physics-informed losses. Leveraged PyTorch for neural network implementation, SciPy for numerical solutions, and Plotly for dynamic visualizations.
\end{highlights}


\begin{twocolentry}{\textit{2024}}
\textbf{Quantitative Finance}
\end{twocolentry}
\begin{highlights}
\item \href{https://github.com/HamidAbbasi-R/Black-Scholes-Model}{\faGithub\hspace{0.13cm}/Black-Scholes-Model}: Created an interactive dashboard explaining the derivation of Black-Scholes model and extensions like implied volatility and American option pricing.
\item \href{https://github.com/HamidAbbasi-R/Stochastic-Processes}{\faGithub\hspace{0.13cm}/Stochastic-Processes}: Developed a dashboard for Portfolio Management: Applied Modern Portfolio Theory to assess the efficient portfolio. Evaluated the Value at Risk (VaR) of each portfolio within a given confidence level.
\item Algorithmic Trading: Designed and backtested different trading strategies based on statistical \href{https://github.com/HamidAbbasi-R/Pair-Arbitrage}{\faGithub\hspace{0.13cm}/Pair-Arbitrage}, range Breakout and retest (\href{https://github.com/HamidAbbasi-R/Adaptive-Channel-Strategy}{\faGithub\hspace{0.13cm}/Adaptive-Channel-Strategy}), and \href{https://github.com/HamidAbbasi-R/Zero-Intelligence-Strategy}{\faGithub\hspace{0.13cm}/Zero-Intelligence-Strategy}
\end{highlights}

\section{Technologies}

\begin{onecolentry}
\textbf{Programming:} \textbf{Python}, \textbf{C++}, \textbf{MATLAB}, \textbf{Fortran}, \textbf{\LaTeX}, \textit{Parallel Computing} (MPI, OpenMP, mpi4py, CUDA, RAPIDS, dask), \textit{High-Dimensional Data Analysis} (Pandas, Polars, Xarray)
\end{onecolentry}

\vspace{0.1 cm}

\begin{onecolentry}
\textbf{Physics Simulation Package}: COMSOL, ANSYS
\end{onecolentry}

% \section{Research Interests}
% % I love sharing knowledge with like-minded people and take pride in being open-minded. Learning fuels me and I have no hesitation in owning up to my mistakes or acknowledging when others have greater expertise. %Beyond that, my specific research interests include:%I’m also passionate about tackling large/complex problems that can’t be solved analytically, which is where approximate models come in.  
% \vspace{0.1 cm}
% \begin{onecolentry}
% \textbf{Numerical Modeling:} Monte Carlo simulations, FVM/FDM for PDEs,  
% \textbf{Machine Learning \& Optimization:} Deep and ensemble learning, Physics-informed ML, Heuristic optimization; 
% \textbf{Quantitative Finance:} Stochastic Calculus, Option pricing, Portfolio Optimization,;
% \textbf{Tomographic Images:} Feature extraction, Visualization

% % \vspace{0.1 cm}


% % \vspace{0.1 cm}


% % \vspace{0.1 cm}


% \end{onecolentry}
% \vspace{0.2 cm}
% \textbf{Risk Management and Portfolio Optimization:} Modern Portfolio Theory
% \end{onecolentry}

% \section{References}
% \begin{onecolentry}
% \item Professor Konstantinos Theodoropoulos \href{mailto:k.theodoropoulos@manchester.ac.uk}{k.theodoropoulos@manchester.ac.uk}
% \item Dr. Masoud Babaei \href{mailto:masoud.babaei@manchester.ac.uk}{masoud.babaei@manchester.ac.uk}
% \end{onecolentry}

\end{document}